\documentclass[11pt]{article}
\usepackage[margin=1in]{geometry}
\usepackage{amsmath,amssymb}
\usepackage{booktabs}
\usepackage{graphicx}
\usepackage{hyperref}
\usepackage{enumitem}
\usepackage[table]{xcolor}
\hypersetup{colorlinks=true,linkcolor=blue,citecolor=blue,urlcolor=blue}
\newcommand{\code}[1]{\texttt{#1}}
\newcommand{\todo}[1]{\textcolor{red}{[TODO: #1]}}

\title{Graduate Deep Reinforcement Learning Paper Chains\\Experimental Report}
\author{DRL Course Repository (projects/grad\_rl)}
\date{\today}

\begin{document}
\maketitle
\begin{abstract}
This report documents reproducible experimental baselines for five canonical paper chains in deep reinforcement learning: value-based, policy-gradient/trust-region, actor--critic with maximum entropy, model-based, and multi-agent. Each section details objectives, algorithms, default hyperparameters, evaluation protocols, and placeholders for results produced by the accompanying Python scripts in \code{projects/grad\_rl/}.
\end{abstract}

\section{How to Reproduce}
\begin{enumerate}[leftmargin=*]
    \item Create and activate a virtual environment, then install dependencies: \code{pip install -r projects/grad\_rl/requirements.txt}.
    \item Run the desired script (examples in Table~\ref{tab:commands}); set \code{--seed} for determinism.
    \item Results are saved under \code{projects/grad\_rl/outputs/<chain>/}. TensorBoard logs live in \code{projects/grad\_rl/outputs/tb/}.
    \item Insert the produced metrics into the result tables below.
\end{enumerate}

\begin{table}[h]
    \centering
    \caption{Entry-point commands (laptop-friendly defaults).}
    \label{tab:commands}
    \begin{tabular}{ll}
        \toprule
        Chain & Command \\
        \midrule
        Value-based & \code{python projects/grad\_rl/scripts/train\_value.py --env ALE/Breakout-v5 --total-steps 200000} \\
        Policy (PPO) & \code{python projects/grad\_rl/scripts/train\_policy.py --algo ppo --env HalfCheetah-v4 --total-steps 300000} \\
        Policy (TRPO) & \code{python projects/grad\_rl/scripts/train\_policy.py --algo trpo --env Walker2d-v4} \\
        Policy (CPO-lite) & \code{python projects/grad\_rl/scripts/train\_policy.py --algo cpo --env SafetyPointCircle0-v0} \\
        Actor--Critic (SAC) & \code{python projects/grad\_rl/scripts/train\_actor\_critic.py --algo sac --env Pendulum-v1} \\
        Model-based (Dyna-Q) & \code{python projects/grad\_rl/scripts/train\_model\_based.py --algo dyna-q} \\
        Model-based (ME-TRPO) & \code{python projects/grad\_rl/scripts/train\_model\_based.py --algo me-trpo --env Pendulum-v1} \\
        Multi-agent (QMIX) & \code{python projects/grad\_rl/scripts/train\_marl.py --algo qmix --iters 50} \\
        Multi-agent (MADDPG) & \code{python projects/grad\_rl/scripts/train\_marl.py --algo maddpg --iters 50} \\
        \bottomrule
    \end{tabular}
\end{table}

\section{Experimental Setup}
\textbf{Hardware}: specify CPU/GPU model and RAM.\\
\textbf{Software}: Python 3.10+, PyTorch 2.2, Gymnasium 0.29, SB3 2.3, Ray 2.9.\\
\textbf{Seeding}: all scripts accept \code{--seed}; vector envs offset seeds per worker.\\
\textbf{Logging}: TensorBoard scalars for reward, loss, KL; JSON metrics saved alongside checkpoints.

\section{Chain 1: Value-Based (DQN, Double, Dueling)}
\subsection{Objective}
Stabilize Q-learning with replay and target networks; reduce maximization bias (Double) and improve representation (Dueling).

\subsection{Default Hyperparameters}
\begin{table}[h]
\centering
\begin{tabular}{ll}
\toprule
Parameter & Value \\
\midrule
Optimizer & Adam, $\alpha=10^{-4}$ \\
Replay buffer & $10^5$ transitions \\
Batch size & 32 \\
Discount $\gamma$ & 0.99 \\
Target update & every 1000 steps \\
Exploration & $\epsilon$-greedy, 0.1 $\to$ 0.01 \\
Dueling & enabled by default \\
\bottomrule
\end{tabular}
\end{table}

\subsection{Metrics}
Mean episode return and std over 5 eval episodes (greedy). Fill Table~\ref{tab:dqnresults} after running.

\begin{table}[h]
\centering
\caption{DQN-family results.}
\label{tab:dqnresults}
\begin{tabular}{lccc}
\toprule
Variant & Steps & Mean Return & Std Return \\
\midrule
DQN & & & \\
Double DQN & & & \\
Dueling DQN & & & \\
QR-DQN (optional) & & & \\
\bottomrule
\end{tabular}
\end{table}

\section{Chain 2: Policy Gradient and Trust Region (REINFORCE, PPO, TRPO, CPO-lite)}
\subsection{Objectives}
\begin{itemize}[leftmargin=*]
    \item REINFORCE: likelihood-ratio gradient with baseline.
    \item TRPO: monotonic improvement via KL-constrained CG step.
    \item PPO: first-order clip surrogate.
    \item CPO-lite: Lagrangian penalty on cost with PPO backbone (requires env exposing \code{info["cost"]}).
\end{itemize}

\subsection{Hyperparameters}
Learning rate $3\times10^{-4}$, GAE $\lambda=0.95$, clip 0.2 (PPO), target KL 0.01 (TRPO), cost limit 0.02 and penalty step 0.05 (CPO-lite).

\subsection{Metrics}
Average return and average cost (if available) over 5 eval episodes. Record in Table~\ref{tab:policyresults}.

\begin{table}[h]
\centering
\caption{Policy-gradient results.}
\label{tab:policyresults}
\begin{tabular}{lcccc}
\toprule
Algo & Env & Steps & Mean Return & Mean Cost \\
\midrule
REINFORCE & & & & \\
PPO & & & & \\
TRPO & & & & \\
CPO-lite & & & & \\
\bottomrule
\end{tabular}
\end{table}

\section{Chain 3: Actor--Critic and Maximum Entropy (A3C/A2C, SAC)}
\subsection{Objectives}
A3C (implemented via synchronous A2C) for on-policy parallel training; SAC for off-policy maximum entropy with temperature tuning.

\subsection{Key Settings}
A2C: $n_{\text{steps}}=5$, $\gamma=0.99$, $\lambda=0.95$. SAC: replay buffer $2\times10^5$, $\tau=0.005$, entropy coefficient auto or 0.2.

\subsection{Metrics}
Mean/Std return over 5 eval episodes (Table~\ref{tab:acresults}).

\begin{table}[h]
\centering
\caption{Actor--critic results.}
\label{tab:acresults}
\begin{tabular}{lccc}
\toprule
Algo & Steps & Mean Return & Std Return \\
\midrule
A2C (A3C stand-in) & & & \\
SAC & & & \\
\bottomrule
\end{tabular}
\end{table}

\section{Chain 4: Model-Based (Dyna-Q, ME-TRPO-style Ensemble)}
\subsection{Objectives}
Dyna-Q mixes real experience with simulated model rollouts for tabular tasks. ME-TRPO trains an ensemble dynamics model and periodically updates a TRPO policy using imagined rollouts to mitigate model bias.

\subsection{Notes}
\begin{itemize}[leftmargin=*]
    \item Dyna-Q uses 20 planning updates per real step; $\epsilon$-greedy with $\epsilon=0.1$.
    \item Ensemble models are two-layer MLPs; horizon for imagination defaults to 5.
    \item The ME-TRPO code is a compact teaching stub; rewards in imagination are not synthesized—treat as a scaffold to extend.
\end{itemize}

\subsection{Metrics}
Report average reward over last 100 episodes (Dyna-Q) and mean eval return (ME-TRPO) in Table~\ref{tab:modelresults}.

\begin{table}[h]
\centering
\caption{Model-based results.}
\label{tab:modelresults}
\begin{tabular}{lccc}
\toprule
Algo & Environment & Metric & Value \\
\midrule
Dyna-Q & CliffWalking-v0 & Avg last-100 return & \\
ME-TRPO & Pendulum-v1 & Mean eval return & \\
\bottomrule
\end{tabular}
\end{table}

\section{Chain 5: Multi-Agent (QMIX, MADDPG)}
\subsection{Objectives}
QMIX factorizes the joint value with a monotonic mixing network for cooperative tasks; MADDPG uses centralized critics and decentralized actors for mixed settings.

\subsection{Settings}
PettingZoo MPE \code{simple\_spread\_v3}, max cycles 25, Ray RLlib workers=2, envs/worker=4, $\gamma=0.99$, $\text{lr}=5\times10^{-4}$.

\subsection{Metrics}
Track episode reward mean per iteration; plot learning curves or record key iterations in Table~\ref{tab:marlresults}.

\begin{table}[h]
\centering
\caption{Multi-agent results.}
\label{tab:marlresults}
\begin{tabular}{lcc}
\toprule
Algo & Iteration & Episode Reward Mean \\
\midrule
QMIX & & \\
MADDPG & & \\
\bottomrule
\end{tabular}
\end{table}

\section{Discussion Template}
\begin{itemize}[leftmargin=*]
    \item \textbf{Stability}: compare variance across seeds for on-policy vs off-policy.
    \item \textbf{Sample Efficiency}: steps to reach target return; note gains from model-based rollouts.
    \item \textbf{Safety}: cost constraint satisfaction in CPO-lite; force violations in Amasa env.
    \item \textbf{Scalability}: wall-clock vs number of workers (A2C/A3C, QMIX).
    \item \textbf{Failure Modes}: overestimation bias (DQN), model bias (ME-TRPO), credit assignment (QMIX).
\end{itemize}

\section{Conclusion}
Summarize the observed trade-offs between stability, efficiency, safety, and scalability across the five chains. Highlight which algorithms suit which environment families and suggest future extensions (e.g., Rainbow, true CPO with conjugate gradient, PlaNet/MuZero for vision, SMACv2 for MARL).

\end{document}
